\documentclass[a4paper, 11pt, UTF8]{article}
%\usepackage{background}
\usepackage{wallpaper}

\usepackage{amssymb,amsmath,amsfonts}
\usepackage{ctex}
\usepackage{setspace}
\usepackage{graphicx}
\usepackage{hyperref}
\usepackage{fontspec}
\usepackage{unicode-math}
\usepackage{ulem}
\usepackage{color}
\usepackage{tabularx}
\usepackage{xargs}
\usepackage[left=1cm,right=1cm,top=1cm,bottom=0cm]{geometry}

\hypersetup{colorlinks=true,linkcolor=blue,urlcolor=cyan}


\parindent 0pt


%\backgroundsetup{scale=1,angle=0,opacity=0.1,contents={\includegraphics[width=\paperwidth, height=\paperheight, keepaspectratio]{b04.jpg}}}

%{ 关于字体设置
\setmainfont{TeXGyrePagella-Regular}
\setmathfont{TeXGyrePagellaMath-Regular}
\setmathfont[range={\mathcal,\Longleftrightarrow,\Longrightarrow,\Longleftarrow,\Rightarrow,\Leftarrow,\Leftrightarrow}]{LatinModernMath-Regular}
\setmathfont[range={\backslash}]{MathJax_Size1-Regular.otf}
\setmathfont[range={\mathbb}]{TeXGyreSchola-Regular}
\setmathfont[range={\complement}]{TeXGyrePagella-Bold}
\setmathfont[range=\mathfrak]{TeXGyrePagella-Regular}
\setmathfont[range={\tilde}]{DejaVuMathTeXGyre.ttf}


\newfontfamily{\tgnr}{TeXGyrePagella-Regular}
\newfontfamily{\tgbf}{TeXGyreSchola-Regular}
\newfontfamily{\tgbfx}{TeXGyrePagella-Bold}
\newfontfamily{\tgbfxx}{FreeSerifBold}
\newfontfamily{\tgsl}{TeXGyrePagella-Italic}
\newfontfamily{\tgsc}{TeXGyrePagella-BoldItalic}
\newfontfamily{\forbra}{MathJax_Size1-Regular.otf}
\newfontfamily{\forbrax}{MathJax_Size2-Regular.otf}
\newfontfamily{\forslash}{MathJax_Size1-Regular.otf}
\def\mySlash{\hspace{-1.2pt}\text{\forslash/}}

\newcommand{\Largenr}[1]{{\Large\tgnr#1}}
\newcommand{\Largesl}[1]{{\Large\tgsl#1}}
\newcommand{\Largebf}[1]{{\Large\tgbf#1}}
\newcommand{\Largebfx}[1]{{\Large\tgbfx#1}}
\newcommand{\Largebfxx}[1]{{\Large\tgbfxx#1}}

\def\envFontHuge{\def\envFontA{\huge}\def\envFont{\LARGE}\def\envFontB{\Large}}
\def\envFontLarge{\def\envFontA{\LARGE}\def\envFont{\Large}\def\envFontB{\large}}
\def\envFontDefault{\def\envFontA{\Large}\def\envFont{\large}\def\envFontB{\normalsize}}
\newcommand{\envFontSmall}[1][\small]{\def\envFontA{\Large}\def\envFont{\normalsize}\def\envFontB{#1}}
\newcommand{\TextE}[1]{\dbsp\envFontLarge\Largesl{#1}\envFontDefault\par\IndentE}
\newcommand{\TextN}[1]{\dbsp\envFontLarge\Largesl{#1}\envFontDefault\par\IndentN}
\newcommand{\TextB}[1]{\dbsp\envFontLarge\Largesl{#1}\envFontDefault\par\IndentB}
\newcommand{\TextNL}[1]{\dbsp\envFontLarge\Largesl{#1}\envFontDefault\par\IndentNL}
%}

%{ 关于排版缩进的便捷指令
%  · {题目内容}\par\IndentB
%  8 {题目内容}\par\IndentN
% 19 {题目内容}\par\IndentNL
% EXAMPLE: {text}\par\IndentE
\def\IndentB{\Blind{\BulletPoint }}
\def\IndentN{\hspace{8.7pt}}
\def\IndentNL{\hspace{16pt}}
\def\IndentE{{\large\Blind{{\BulletPoint} {\Example}\,\,\,}}}

% 用于对齐(a)(b)(c)(d)
%(a) {text} \par\quad\Ha
%... \par\quad\Ha
%{text} \par\quad
%(b) {text} \par\quad\Hb
%... \par\quad\Ha
%{text} \par\quad
%(c) {text} \par\quad\Hc
%... \par\quad\Ha
%{text} \par\quad
%(d) {text} \par\quad\Hd
%... \par\quad\Ha
%{text} \par\quad
\def\Ha{{\large\Blind{(a) }}}
\def\Hb{{\large\Blind{(b) }}}
\def\Hc{{\large\Blind{(c) }}}
\def\Hd{{\large\Blind{(d) }}}

% 用于对齐(i)(ii)(iii)
%(i) {text} \par\quad\Hi
%... \par\quad\Hi
%{text} \par\quad\Endi
%(ii)  {text} \par\quad\Hii
%... \par\quad\Hii
%{text} \par\quad\Endii
%(iii) {text} \par\quad\Hiii
%... \par\quad\Hiii
%{\text}\par\quad
\def\Endi{}
\def\Endii{}
\def\Hi{\Blind{(i) }}
\def\Hii{\Blind{(ii) }}
\def\Hiii{}

% 用于对齐(I)(II)(III)
\def\EndI{}
\def\EndII{}
\def\EndIII{}
\def\HI{\Blind{(I) }}
\def\HII{\Blind{(II) }}
\def\HIII{\Blind{(III) }}
%}

%{ 关于文本环境
\newcommand{\hMath}[5][-4pt]{#3\hspace{#1}\begin{array}{#2}#5\end{array}\hspace{#1}#4}
\newcommand{\MathLeftBrace}				[2]{\hMath[0pt]{#1}{\left\{}{\right.}{#2}}
\newcommand{\MathRightBrace}			[2]{\hMath[0pt]{#1}{\left.}{\right\}}{#2}}
\newcommand{\MathLeftrightBrace}		[2]{\hMath[0pt]{#1}{\left\{}{\right\}}{#2}}
\newcommand{\MathLeftMid}				[2]{\hMath[0pt]{#1}{\left|}{\right.}{#2}}
\newcommand{\MathRightMid}				[2]{\hMath[0pt]{#1}{\left.}{\right|}{#2}}
\newcommand{\MathLeftrightMid}			[2]{\hMath[0pt]{#1}{\left|}{\right|}{#2}}
\newcommand{\MathLeftrightPare}			[2]{\hMath[0pt]{#1}{\left(}{\right)}{#2}}

\newcommand{\hText}[4][-4pt]{\hMath[#1]{c}{#2}{#3}{#4}}
\newcommand{\Par}[2][1pt]{{\text{\forbra\envFontB(\hspace{#1}}}#2{\text{\hspace{#1}\forbra\envFontB)}}}%{\hText[-4pt]{}{}{$\hspace{-1pt}{\envFontB$\hText[#1]{\left(}{\right)}{${\envFont$#2$}$}$}$}}
\newcommandx{\Mid}[3][1=1pt,2=\envFontA]{{\text{#2\mid}}\hspace{#1}{#3}\hspace{#1}{\text{#2\mid}}}%{\hText[-2pt]{}{}{$\hspace{-1pt}{\envFontB$\hText[#1]{\left|}{\right|}{${\envFont$#2$}$}$}$}}
\newcommand{\cMid}[2][1pt]{\Mid[#1][\envFontLarge\envFontA]{#2}}
\newcommand{\aMid}[2][1pt]{\Mid[#1][\envFontLarge\envFontA]{#2}}
\newcommand{\aXMid}[2][1pt]{\Mid[#1][\envFontHuge\envFont]{#2}}
\newcommand{\Bra}[2][-5pt]{{\text{\forbra\envFont{\{}}}#2{\text{\forbra\envFont\}}}}%{\hText[-4pt]{}{}{$\hspace{-1pt}{\envFontB$\hText[#1]{\left\{}{\right\}}{${\envFont$#2$}$}$}$}}
\newcommand{\SmallPar}[2][0pt]{{\envFontSmall[\scriptsize]\Par[#1]{#2}}}
\newcommand{\BigPar}[1]{{\text{\envFont\forbra(\hspace{1pt}}}#1{\text{\hspace{1pt}\envFont\forbra)}}}
\newcommand{\XPar}[1]{{\text{\envFont\forbrax(\hspace{1pt}}}#1{\text{\hspace{1pt}\envFont\forbrax)}}}
\newcommand{\Sbra}[1]{{\text{\envFont\forbra[\hspace{1pt}}}#1{\text{\envFont\forbra\hspace{1pt}]}}}
\newcommand{\XSbra}[1]{{\text{\envFont\forbrax[}}#1{\text{\envFont\forbrax]}}}
\newcommand{\zeroSubs}[1][\Bra]{{#1{0}}}

\newcommand{\Interval}[3]{{\text{\envFont\forbra#1\hspace{1pt}}}#3{\text{\envFont\forbra\hspace{1pt}#2}}}
\newcommand{\XInterval}[3]{{\text{\envFont\forbrax#1}}#3{\text{\envFont\forbrax#2}}}

\newcommand{\ExampleX}[2][\Solution]{
	{\BulletPoint} {\Example}\,\,\,{#2}\envFontDefault\vspace{-16pt}\par
	#1
}
\newcommand{\ProblemN}[3][\Solution]{
	{\Onumber{#2}} {#3}\envFontDefault\vspace{-16pt}\par
	#1
}
\newcommandx{\ProblemNor}[5][1=\Solution,4=\Sbra]{
	{\Onumber{#2}} \;{\dbsp#4{\OR({\normalsize#3})}}{#5}\envFontDefault\vspace{-16pt}\par
	#1
}
\newcommandx{\ProblemNnoor}[5][1=\Solution,4=\Sbra]{
	{\Onumber{#2}} \;{\dbsp#4{\small#3}}{#5}\envFontDefault\vspace{-16pt}\par
	#1
}
\newcommand{\ProblemB}[2][\Solution]{
	{\BulletPoint}{#2}\envFontDefault\vspace{-16pt}\par
	#1
}
\newcommand{\ProblemBc}[2][\Solution]{
	{$\circ$}{#2}\envFontDefault\vspace{-16pt}\par
	#1
}
\newcommandx{\ProblemBnoor}[4][1=\Solution,3=\Par]{
	{\BulletPoint} \dbsp#3{{\small#2}}{#4}\envFontDefault\vspace{-16pt}\par
	#1
}
\newcommand{\ProblemBor}[3][\Solution]{
	{\BulletPoint} {\dbsp\OR({\normalsize#2})}{#3}\envFontDefault\vspace{-16pt}\par
	#1
}
\newcommand{\AlignEq}[2]{
	\vspace{-25pt}
	\begin{align*}
		#1#2
	\end{align*}
	\vspace{-25pt}
}
%}

%{ 关于常用标识的便捷指令

\newcommand{\Frac}[2]{{\displaystyle{\frac{#1}{#2}}}}
\def\notRightarrow{\Rightarrow\!\!\!\!\!\!\!/\,\,\,\,}%\hspace{-13.5pt}/\hspace{7.5pt}}
\def\notLeftarrow{\Leftarrow\hspace{-13.5pt}/\hspace{7.5pt}}
\def\Nbb{{\mathbb{N}}}
\def\Zbb{{\mathbb{Z}}}
\def\Qbb{{\mathbb{Q}}}
\def\Fbb{{\mathbb{F}}}
\def\Rbb{{\mathbb{R}}}
\def\Cbb{{\mathbb{C}}}
\def\Nbp{{\mathbb{N}^+}}
\def\mmid{\text{\envFontA|}}
\def\d{{\textup{\tgnr d}}}
\def\i{{\textup{\tgnr i}}}
%\def\dim{{\textup{\tgnr dim}}\,}
\def\Dim{{\textup{\tgnr dim}}}
\def\Deg{{\textup{\tgnr deg}}}
\def\range{{\textup{\tgnr range}}\,}
\def\null{{\textup{\tgnr null}}\,}
\def\card{{\textup{\tgnr card}}\,}
\def\Spn{{\textup{\tgnr span}}\,}
\newcommand{\Span}[2][\Par]{{\textup{\tgnr span}}#1{#2}}
\newcommand{\Lm}[2][\Par]{\mathcal{L}#1{#2}}
\newcommand{\Lfm}[2][\Par]{#1{#2\apostrophe}}
\newcommand{\Mt}[2][\Par]{\mathcal{M}#1{#2}}
\def\Mneg{\mathcal{M}^{-1}}
\def\Po{\mathcal{P}}
\newcommand{\PoF}[2][\Par]{\Po_{#2}#1{\Fbb}}
\newcommand{\PoR}[2][\Par]{\Po_{#2}#1{\Rbb}}
\newcommand{\PoC}[2][\Par]{\Po_{#2}#1{\Cbb}}
\def\mathC{C}
\def\apostrophe{\prime}
\def\BulletPoint{{\small\bullet}}
\def\BulletPointX{\BulletPoint \,\hspace{1pt}}
\def\bullpt{{\tiny\bullet}}

\def\OR{{\large O{\footnotesize R} }}
\def\Or{{\large O{\footnotesize R.} }}
\def\Solution{{\tgbfx\large S\footnotesize{OLUTION:}}\,\,\,}
\def\NOTE{\tgnr\large N{\footnotesize OTE}}
\def\NOTEFOR{{\tgnr\large N{\footnotesize OTE} F{\footnotesize OR}}}
\def\NEWTHEOREM{{\tgnr\large N{\footnotesize EW} T{\footnotesize HEOREM}}}
\def\NOTICE{{\tgnr\large N{\footnotesize OTICE}\;}}
\def\COMMENT{{\tgnr\large C{\footnotesize OMMENT}\;}}
\def\COROLLARY{{\tgnr\large C{\footnotesize OROLLARY}}}
\def\TIPS{{\tgnr\large T{\footnotesize IPS}\;}}
\def\Tips{{\tgbfx\large T{\footnotesize IPS}:}}
\newcommand{\TipsN}[1]{{\tgbfx\large T{\footnotesize IPS {\large#1\,}}:}}
\newcommand{\Onumber}[1]{\Largebfxx{#1}\hspace{-2pt}}
\newcommand{\NoteFor}[1]{\Largebfx{N{\small OTE} F{\small OR} #1:}}
\newcommand{\NoteForSmall}[1]{{\tgbfx\large N{\footnotesize OTE} F{\footnotesize OR} #1:}}
\newcommand{\NewNotation}{{\tgbfx\large N{\footnotesize EW} N{\footnotesize OTATION}:}}
\newcommand{\NewTheorem}{{\tgbfx\large N{\footnotesize EW} T{\footnotesize HEOREM}:}}
\newcommand{\Comment}{{\tgbfx\large C\footnotesize{OMMENT:}}}
\newcommand{\Example}{{\tgbfx\large E\footnotesize{XAMPLE:}}}
\newcommand{\Exercise}[1]{{\tgbfx\large E{\footnotesize XERCISE} {#1}\hspace{2pt}:}}
\newcommand{\Corollary}{{\tgbfx\large C{\footnotesize OROLLARY}:}}

\def\ChEnd{\rightline{\Largebfx{E{\small NDED}}}\par\vspace{6pt}}
\newcommand{\PfEnd}[1][-18pt]{{\large\vspace{#1}\par\hfill$\square$\par}}
\newcommandx{\SepLine}[2][1=0pt,2={}]{{\vspace{-5pt}\par
		#2\tiny \_\,\_\,\_\,\_\,\_\,\_\,\_\,\_\,\_\,\_\,\_\,\_\,\_\,\_\,\_\,\_\,\_\,\_\,\_\,\_\,\_\,\_\,\_\,\_\,\_\,\_\,\_\,\_\,\_\,\_\,\_\,\_\,\_\,\_\,\_\,\_\,\_\,\_\,\_\,\_\,\_\,\_\,\_\,\_\,\_\,\_\,\_\,\_\,\_\,\_\,\_\,\_\,\_\,\_\,\_\,\_\,\_\,\_\,\_\,\_\,\_\,\_\,\_\,\_\,\_\,\_\,\_\,\_\,\_\,\_\,\_\_\,\_\,\_\,\_\,\_\,\_\,\_\,\_\,\_\,\_\,\_\,\_\,\_\,\_\,\_\,\_\,\_\,\_\,\_\,\_\,\_\,\_\,\_\,\_\,\_\,\_\,\_\,\_\,\_\,\_\,\_\,\_\,\_\,\_\,\_\,\_\,\_\,\_\,\_\,\_\,\_\,\_\,\_\,\_\,\_\,\_\,\_\,\_\,\_\,\_\,\_\,\_\,\_\,\_\,\_\,\_\,\_\,\_\,\_\,\_\,\_\,\_\,\_\,\_\,\_\,\_\,\_\,\_\,\_\,\_\,\_}\vspace{5pt}\vspace{#1}\par}
\newcommand{\ChDecl}[3]{{\huge\tgbfxx\hypertarget{#1}{#2}}{#3}{\vspace{5pt}}}
\newcommand\hLk[2]{\hyperlink{#1}{#2}}
\newcommand\Lch[2]{\hLk{Ch#1}{#2}}
\newcommand\TXT[1]{\textup{#1}}
\newcommand{\Blind}[1]{\textcolor{white}{#1}}
%}

%\orMode[setup]{pdf version}[space]{print version}

%{ 关于调试、版本切换
\def\dbsp{\bullet}
\def\dbsp{} % 方便调试段落缩进
\newcommandx{\orMode}[5][1=\hMath{l}{}{},2=\hMath{l}{}{},4=-18pt]{{\normalsize
% If you want the debug version.
	%{\qquad$#1{$#3$}$}\vspace{#4}\par{\qquad$#2{$#5$}$}
% If you want the pdf version.
	{\qquad$#1{$#3$}$}
% If you want the print version.
	%{\qquad$#2{$#5$}$}
}}
%}

